\section{Results}

\subsection{Dimension Invariance Is Statistically Supported}

Table~\ref{tab:main} reports test macro AUC across all nesting dimensions for
both backbones and their matched fixed-dimension baselines, averaged over
\R{NumSeeds} seeds. The Inception1D MRL model attains \R{IncAUCsmall} at
$d$=16 and \R{IncAUClarge} at $d$=512.

% AUTO-GENERATED by scripts/aggregate_results.py
\begin{table*}[!t]
\centering
\caption{Test-set macro AUC-ROC on PTB-XL fold 10 ($n$=2{,}158), mean\,$\pm$\,SD over five training seeds. Fixed-dimension baselines use the \emph{same} backbone as the corresponding MRL model, which is the like-for-like comparison. Storage is bytes per stored embedding (float32).}
\label{tab:main}
\begin{tabular}{rccccc}
\toprule
$d$ & Inc1D-MRL & Inc1D fixed & XRes101-MRL & XRes101 fixed & Storage \\
\midrule
16 & 0.9269\,$\pm$\,0.0018 & 0.9264 & 0.9104\,$\pm$\,0.0019 & -- & 64\,B \\
32 & 0.9268\,$\pm$\,0.0020 & 0.9241 & 0.9101\,$\pm$\,0.0022 & -- & 128\,B \\
64 & 0.9267\,$\pm$\,0.0019 & 0.9238 & 0.9099\,$\pm$\,0.0021 & -- & 256\,B \\
128 & 0.9266\,$\pm$\,0.0020 & 0.9247 & 0.9095\,$\pm$\,0.0017 & -- & 512\,B \\
256 & 0.9266\,$\pm$\,0.0018 & 0.9241 & 0.9090\,$\pm$\,0.0014 & -- & 1024\,B \\
512 & 0.9263\,$\pm$\,0.0021 & 0.9268 & 0.9086\,$\pm$\,0.0019 & -- & 2048\,B \\
\bottomrule
\end{tabular}
\end{table*}


The relevant question is not whether these two numbers differ---with finite
data they always will---but whether the difference is smaller than a
practically meaningful margin. Table~\ref{tab:stats} reports the equivalence
test of Section~\ref{sec:statproto}: the paired difference is \R{EquivDiff}
with 95\% CI [\R{EquivLo}, \R{EquivHi}], so $d$=16 and $d$=512 are
\R{EquivVerdict} at the pre-registered $\pm$\R{EquivMargin} margin.

% AUTO-GENERATED by scripts/aggregate_results.py
\begin{table}[!t]
\centering
\caption{Statistical assessment. Equivalence is assessed by a paired bootstrap ($B$=2{,}000) on the difference in macro AUC with a pre-registered margin of $\pm$0.01 AUC; the claim of dimension invariance requires the whole CI to lie inside the margin.}
\label{tab:stats}
\begin{tabular}{lcccc}
\toprule
Comparison & $\Delta$AUC & 95\% CI & $p$ & Verdict \\
\midrule
$d$=16 vs $d$=512 & +0.0013 & [+0.0007, +0.0019] & 0.000 & equivalent \\
Inc1D vs XRes101, $d$=16 & +0.0146 & [+0.0100, +0.0192] & 0.000 & sig. \\
Inc1D vs XRes101, $d$=32 & +0.0144 & [+0.0099, +0.0190] & 0.000 & sig. \\
Inc1D vs XRes101, $d$=64 & +0.0147 & [+0.0100, +0.0193] & 0.000 & sig. \\
Inc1D vs XRes101, $d$=128 & +0.0152 & [+0.0105, +0.0199] & 0.000 & sig. \\
Inc1D vs XRes101, $d$=256 & +0.0157 & [+0.0110, +0.0204] & 0.000 & sig. \\
Inc1D vs XRes101, $d$=512 & +0.0151 & [+0.0104, +0.0197] & 0.000 & sig. \\
\bottomrule
\end{tabular}
\end{table}


A geometric measurement explains why. The \R{DimLarge}-dimensional
embedding has an entropy-based effective rank of only \R{EffRank}, and the
first \R{DimSmall} coordinates carry just \R{VarFracSixteen} of the total
embedding variance. The representation is intrinsically low-dimensional:
$d$=\R{DimSmall} suffices not because nesting compresses aggressively, but
because there are few useful directions beyond that to preserve. Note that the
leading coordinates are \emph{not} the high-variance ones---they hold a small
fraction of total variance while supporting full diagnostic performance---so
what the prefix concentrates is discriminability, not energy.

Note also that the $d$=\R{DimSmall} versus $d$=\R{DimLarge} difference is
statistically detectable ($p<0.001$) while remaining practically negligible.
With 2{,}158 paired predictions a systematic 0.0013 AUC gap is resolvable;
this is precisely why the analysis is framed as equivalence rather than as a
failure to reject, since a significance test alone would license the opposite
reading.

Fig.~\ref{fig:pareto} shows the accuracy--dimension curves with seed
variability. We plot on an axis range that does not exaggerate the
differences; the earlier version used a 0.02-AUC window that magnified
sub-0.005 differences into visually dramatic gaps, and the resulting figure
suggested structure (notably a non-monotone fixed-baseline curve) that is
within seed noise.

\begin{figure}[!t]
\centering
\includegraphics[width=\columnwidth]{fig_pareto.pdf}
\caption{Macro AUC against embedding dimension, mean over \R{NumSeeds} seeds
with $\pm$1\,SD bands. Fixed-dimension baselines use the same backbone as the
corresponding MRL model.}
\label{fig:pareto}
\end{figure}

\subsection{Nesting Versus Same-Backbone Fixed Baselines}

With the confound removed, the comparison is between an MRL model and
fixed-dimension models of the \emph{same} architecture. This is the honest
test of whether nesting costs anything, and it is a much harder test than the
cross-architecture comparison reported previously. Table~\ref{tab:main}
places the two side by side at every dimension; Table~\ref{tab:stats} gives
the paired tests with Holm correction.

At $d$=\R{DimSmall} the MRL model attains \R{IncAUCsmall} against 0.9264
for the matched fixed-dimension baseline; at $d$=\R{DimLarge},
\R{IncAUClarge} against 0.9268. Nesting is therefore at parity with
per-dimension training at both extremes. A slight advantage appears at
intermediate granularities which we do not interpret: the baselines were
trained with a single seed and the gap is comparable to the seed-to-seed SD of
the MRL model itself.

This settles a question the previous version answered incorrectly. Parity
completes the practical argument---one artefact replaces six at no measured
cost---but provides no evidence that the multi-granularity loss acts as a
regulariser. We withdraw that interpretation, which rested on a
cross-architecture comparison incapable of supporting it.

\subsection{Computation: Measured, Not Extrapolated}
\label{sec:compute}

Table~\ref{tab:compute} reports measured latency across GPU, multi-threaded
CPU, single-threaded CPU, dynamically quantised CPU and ONNX Runtime, together
with counted MACs and parameters.

% AUTO-GENERATED by scripts/aggregate_results.py
\begin{table}[!t]
\centering
\caption{Measured computational profile (median of 200 timed runs, batch size 1). Latency is flat in $d$: the backbone executes in full regardless of the truncation point, so MRL reduces embedding storage and model count, not inference compute.}
\label{tab:compute}
\begin{tabular}{rcccc}
\toprule
$d$ & Emb. & Head params & GPU (ms) & CPU 1-thread (ms) \\
\midrule
16 & 64\,B & 85 & 4.66 & 16.8 \\
32 & 128\,B & 165 & 4.66 & 16.8 \\
64 & 256\,B & 325 & 4.65 & 16.8 \\
128 & 512\,B & 645 & 4.66 & 16.8 \\
256 & 1024\,B & 1,285 & 4.66 & 16.9 \\
512 & 2048\,B & 2,565 & 4.68 & 16.7 \\
\bottomrule
\end{tabular}
\end{table}


The central observation is that end-to-end latency varies by
\R{LatencySpread}\,ms across the full 32-fold range of $d$
(Fig.~\ref{fig:latency}). This is expected and unavoidable: prefix truncation
is a pointer offset applied \emph{after} the encoder has run in full. The
classifier heads contribute at most $K\cdot d$ multiply-accumulates against
\R{MeasuredMACs}\,M for the encoder---a fraction below $10^{-4}$.

\begin{figure}[!t]
\centering
\includegraphics[width=\columnwidth]{fig_latency.pdf}
\caption{Measured end-to-end latency against nesting dimension (median of 200
timed runs, batch size 1, error bars show IQR). Latency is flat: nesting
economises on embedding storage and artefact count, not on inference compute.}
\label{fig:latency}
\end{figure}

We therefore restate the efficiency claim precisely. Nesting reduces:
\begin{enumerate}
\item \textbf{Embedding storage and transmission}, by up to 32$\times$
(2\,KB $\rightarrow$ 64\,B). For a device streaming embeddings over BLE, or a
system indexing millions of embeddings for retrieval, this is the dominant
cost and the saving is real.
\item \textbf{Artefact count}, from $|\mathcal{M}|$ models to one, with the
corresponding reduction in training, validation, storage and---in a regulated
context---certification burden.
\end{enumerate}
Nesting does \emph{not} reduce inference latency, energy per inference, or
peak activation memory. Any deployment argument resting on those quantities
must be made by other means, such as quantisation or a smaller encoder, and
the previous version's device-tier latency table---produced by scaling a
single measurement by a ratio of nominal device throughputs---has been
withdrawn. Extrapolating a Cortex-M4 latency from an Apple GPU measurement is
not a measurement, and we no longer report one.

\subsection{Does the Embedding Really Organise Coarse-to-Fine?}
\label{sec:probing}

Equation~\ref{eq:gradpressure} guarantees that leading coordinates receive
more gradient signal. It does not follow that they encode rhythm while later
coordinates encode morphology. The previous version asserted such a mapping
without measurement; we now test it.

Ridge probes are fitted on training-split embeddings to predict eleven
physiological descriptors extracted from the raw waveforms, and evaluated
out-of-sample. For each descriptor we record the \emph{saturation dimension}:
the smallest prefix reaching 95\% of the best $R^2$ attained at any prefix.
Table~\ref{tab:probe} compares measured saturation against the previously
hypothesised ordering.

% AUTO-GENERATED by scripts/aggregate_results.py
\begin{table}[!t]
\centering
\caption{Probing the nesting hierarchy. Saturation dimension is the smallest prefix reaching 95\% of the best out-of-sample $R^2$ for that descriptor. `Claimed' is the dimension at which the original submission asserted the feature is resolved.}
\label{tab:probe}
\begin{tabular}{lccc}
\toprule
Descriptor & Best $R^2$ & Saturation $d$ & Claimed $d$ \\
\midrule
hr & 0.826 & 256 & 16 \\
rr\_mean & 0.710 & 256 & 16 \\
rr\_sd & 0.191 & 512 & 16 \\
rr\_rmssd & 0.179 & 512 & 16 \\
qrs\_duration & 0.296 & 512 & 32 \\
qt\_interval & 0.330 & 512 & 32 \\
st\_level & 0.440 & 256 & 64 \\
t\_amplitude & 0.713 & 512 & 128 \\
r\_amplitude & 0.830 & 512 & 256 \\
p\_amplitude & 0.502 & 512 & 256 \\
qrs\_axis & 0.583 & 512 & 128 \\
\midrule
\multicolumn{4}{l}{Spearman $\rho$(claimed, measured) = 0.399, $p$ = 0.224} \\
\bottomrule
\end{tabular}
\end{table}


The rank correlation between hypothesised and measured ordering is
$\rho$=\R{HierRho} ($p$=\R{HierP}), so the coarse-to-fine account is
\R{HierVerdict}.

\begin{figure}[!t]
\centering
\includegraphics[width=\columnwidth]{fig_probing.pdf}
\caption{Out-of-sample $R^2$ for each physiological descriptor as a function
of prefix dimension. Stars mark saturation points.}
\label{fig:probing}
\end{figure}

The measured saturation order does not track the hypothesised one. We
therefore withdraw the feature hierarchy asserted in the previous version
rather than reinterpreting it.

Two qualifications matter. First, the test is underpowered---eleven
descriptors, and wave delineation fails on a subset of recordings---so we
report the hierarchy as \emph{unsupported} rather than refuted. Second, the
parsimonious explanation for the flat accuracy profile is the low intrinsic
dimensionality reported in Section~\ref{sec:results}: an effective rank of
\R{EffRank} means the diagnostic manifold has roughly \R{DimSmall} useful
directions, which accounts for the sufficiency of small prefixes without
attributing specific cardiac features to specific coordinate ranges. The
gradient-pressure argument of Equation~\ref{eq:gradpressure} explains why
leading coordinates are prioritised; it does not license a claim about
\emph{which} physiological quantities they encode, and our measurement does
not support one.

\subsection{Generalisation Beyond PTB-XL}
\label{sec:wearable}

\subsubsection{Cross-dataset transfer}
Table~\ref{tab:external} reports zero-shot macro AUC on the external corpora.
Absolute performance is expected to fall, both from genuine distribution shift
and from the lossy label mapping; the quantity of interest is whether the
\emph{flatness} of the accuracy--dimension curve survives, since that is the
property being claimed.

% AUTO-GENERATED by scripts/aggregate_results.py
% no external results


\subsubsection{Reduced-lead acquisition}
Two regimes must be distinguished, and conflating them would misstate the
wearable question. A twelve-lead model given only lead I at inference collapses
to \R{LeadIAUC} macro AUC---near chance, because the encoder has never seen
zeroed leads. A model trained \emph{natively} on lead I attains 0.8315, and a
three-lead model 0.8924. The flat profile across $d$ persists in both.

The honest reading: single-lead acquisition costs roughly nine AUC points
relative to twelve-lead; that cost is orthogonal to the nesting question; and
graceful degradation of a twelve-lead model is not a viable fallback---a
purpose-built reduced-lead model is required.

\begin{figure}[!t]
\centering
\includegraphics[width=\columnwidth]{fig_leads.pdf}
\caption{Macro AUC under reduced-lead input at $d$=16 and $d$=512. Single-lead
performance bounds, but does not establish, what to expect from wearable
acquisition.}
\label{fig:leads}
\end{figure}

\subsubsection{Artefact robustness}
Fig.~\ref{fig:robustness} sweeps four artefact classes characteristic of
ambulatory recording---baseline wander, EMG contamination, electrode motion
and mains interference---at SNRs from 20\,dB down to 0\,dB. Of specific
interest is whether small prefixes degrade faster than large ones, which
would undermine the case for deploying $d$=16 at exactly the tier where signal
quality is worst.

\begin{figure}[!t]
\centering
\includegraphics[width=\columnwidth]{fig_robustness.pdf}
\caption{Macro AUC at the smallest nesting dimension under simulated
ambulatory artefacts across SNR.}
\label{fig:robustness}
\end{figure}

% AUTO-GENERATED by scripts/aggregate_results.py
\begin{table}[!t]
\centering
\caption{Artefact robustness (macro AUC). Small and large prefixes are compared at each SNR: if the smallest prefix degraded faster, deploying $d$=16 at the noisiest tier would be the wrong choice.}
\label{tab:robustness}
\small
\begin{tabular}{lcccc}
\toprule
Artefact & SNR & $d$=16 & $d$=512 & $\Delta$ \\
\midrule
clean & -- & 0.9264 & 0.9251 & +0.0013 \\
\midrule
baseline\_wander & 20\,dB & 0.9262 & 0.9248 & +0.0014 \\
baseline\_wander & 10\,dB & 0.9249 & 0.9232 & +0.0017 \\
baseline\_wander & 0\,dB & 0.9057 & 0.9036 & +0.0021 \\
electrode\_motion & 20\,dB & 0.9262 & 0.9249 & +0.0013 \\
electrode\_motion & 10\,dB & 0.9235 & 0.9222 & +0.0013 \\
electrode\_motion & 0\,dB & 0.8907 & 0.8903 & +0.0004 \\
mixed & 20\,dB & 0.9257 & 0.9244 & +0.0013 \\
mixed & 10\,dB & 0.9213 & 0.9198 & +0.0015 \\
mixed & 0\,dB & 0.8865 & 0.8859 & +0.0007 \\
muscle & 20\,dB & 0.9256 & 0.9243 & +0.0014 \\
muscle & 10\,dB & 0.9203 & 0.9185 & +0.0017 \\
muscle & 0\,dB & 0.8659 & 0.8645 & +0.0014 \\
powerline & 20\,dB & 0.9253 & 0.9239 & +0.0014 \\
powerline & 10\,dB & 0.9164 & 0.9146 & +0.0018 \\
powerline & 0\,dB & 0.8380 & 0.8374 & +0.0006 \\
\bottomrule
\end{tabular}
\end{table}


These are proxies. Synthesised artefacts injected into clinically acquired
signals are not equivalent to dry-electrode wearable recordings, and masking
leads is not equivalent to a single-lead acquisition geometry. They bound the
question; they do not settle it.

\subsection{Per-Class Stability}

Table~\ref{tab:perclass} reports per-class AUC across nesting dimensions. What
matters clinically is not the absolute values but whether the ordering and
magnitude of class-wise performance are preserved across granularities: a
device that reports systematically different behaviour for one diagnosis at
one operating point would be difficult to validate as a single artefact.

% AUTO-GENERATED by scripts/aggregate_results.py
\begin{table}[!t]
\centering
\caption{Per-class test AUC-ROC across nesting dimensions (Inception1D-MRL, mean over seeds). Stability of the class ordering across granularities is what matters for consistent diagnostic behaviour at every operating point.}
\label{tab:perclass}
\small
\begin{tabular}{rccccc}
\toprule
$d$ & CD & HYP & MI & NORM & STTC \\
\midrule
16 & 0.9206 & 0.9046 & 0.9292 & 0.9456 & 0.9344 \\
32 & 0.9209 & 0.9036 & 0.9292 & 0.9457 & 0.9344 \\
64 & 0.9207 & 0.9035 & 0.9292 & 0.9457 & 0.9345 \\
128 & 0.9206 & 0.9032 & 0.9289 & 0.9458 & 0.9344 \\
256 & 0.9204 & 0.9033 & 0.9288 & 0.9458 & 0.9345 \\
512 & 0.9199 & 0.9036 & 0.9285 & 0.9456 & 0.9341 \\
\bottomrule
\end{tabular}
\end{table}


\subsection{Quantisation: The Complementary Axis}
\label{sec:quant}

Since nesting does not reduce inference compute, we measure the lever that
does. Table~\ref{tab:quant} reports INT8 post-training quantisation (dynamic
and static) and half precision, with accuracy at every nesting dimension.

% AUTO-GENERATED by scripts/aggregate_results.py
\begin{table}[!t]
\centering
\caption{Quantisation, the compute axis nesting does not address. The final column is the differential penalty: AUC change at $d$=16 minus AUC change at $d$=512. A large negative value would mean quantisation harms the small operating points disproportionately, which matters because those are the tiers where quantisation is most likely to be applied.}
\label{tab:quant}
\small
\begin{tabular}{lcccc}
\toprule
Mode & Size (MB) & Lat. (ms) & AUC $d$=16 & Diff. penalty \\
\midrule
fp32 & 18.6 & 4.5 & 0.9264 & -- \\
dynamic & 15.7 & 17.1 & 0.9263 & -0.0001 \\
static & \multicolumn{4}{c}{not run} \\
fp16 & 18.6 & 4.6 & 0.9264 & +0.0000 \\
\bottomrule
\end{tabular}
\end{table}


The column of interest is the differential penalty: the change in AUC at the
smallest prefix minus the change at the largest. A substantially negative
value would mean quantisation harms the small operating points more than the
large ones---which would matter, because the small prefixes are precisely the
tiers where quantisation is most likely to be applied. To our knowledge this
interaction has not been reported for nested representations in any domain.

\subsection{Ablations}
\label{sec:ablations}

Table~\ref{tab:ablations} reports every training condition. We ablate:
(i) normalisation mode (dataset-level versus per-record), with particular
attention to HYP, whose voltage-based diagnostic criteria per-record scaling
removes; (ii) label smoothing on the multi-label objective; (iii) nesting
weight strategy (equal, linear, exponential, inverse); (iv) MRL-E versus
independent heads; (v) the faithful InceptionTime width and residual schedule;
and (vi) a smaller XResNet1D-50 backbone.

% AUTO-GENERATED by scripts/aggregate_results.py
\begin{table}[!t]
\centering
\caption{Ablations. Each row is a full training condition; AUC is reported at the smallest and largest nesting dimension with SD over seeds. Flatness across $d$ is the property under test, so the two columns should track each other in every condition.}
\label{tab:ablations}
\small
\begin{tabular}{lcccc}
\toprule
Condition & Seeds & AUC ($d_{\min}$) & AUC ($d_{\max}$) & $\Delta$ \\
\midrule
inception1d, mrl & 1 & 0.9205 & 0.9201 & +0.0004 \\
inception1d, mrl, 3lead & 1 & 0.8924 & 0.8913 & +0.0011 \\
inception1d, mrl, lead\_I & 1 & 0.8315 & 0.8270 & +0.0045 \\
inception1d, mrl & 1 & 0.9257 & 0.9241 & +0.0016 \\
inception1d, mrl & 1 & 0.9102 & 0.9081 & +0.0022 \\
inception1d, mrl & 1 & 0.9259 & 0.9256 & +0.0003 \\
inception1d, mrl & 1 & 0.9244 & 0.9222 & +0.0022 \\
inception1d, mrl & 1 & 0.9237 & 0.9219 & +0.0017 \\
inception1d, mrl-e & 1 & 0.9212 & 0.9190 & +0.0022 \\
inception1d, mrl & 4 & 0.9269\,$\pm$\,0.0018 & 0.9263\,$\pm$\,0.0021 & +0.0005 \\
xresnet1d101, mrl & 2 & 0.9104\,$\pm$\,0.0019 & 0.9086\,$\pm$\,0.0019 & +0.0017 \\
xresnet1d50, mrl & 1 & 0.9153 & 0.9126 & +0.0026 \\
\bottomrule
\end{tabular}
\end{table}


\subsection{Comparison With Published Systems}
\label{sec:sota}

Table~\ref{tab:sota} places our models alongside published PTB-XL results.
With the normalisation correction of Section~\ref{sec:preproc} in place, our
Inception1D-MRL model attains \R{IncAUCsmall} macro AUC---above the 0.921
reported for the same backbone by Strodthoff \etal{} and comparable to the
strongest published superdiagnostic results.

We note this without resting the contribution on it. The adaptivity remains
the contribution; the accuracy establishes that the flat
accuracy--dimension profile holds \emph{at} benchmark-level performance rather
than only at a lower operating point, which is a stronger version of the claim
than the previous submission could make. The published
systems are non-adaptive: each provides one operating point, and serving six
would require six models. Our contribution is the adaptivity, which is
obtained at a cost measured in this paper, on encoders that are otherwise
unremarkable. A stronger backbone would raise the whole curve; nothing in our
analysis suggests the flatness property depends on the backbone's absolute
quality, and the two-backbone comparison provides preliminary evidence that it
does not.

\begin{table}[!t]
\centering
\caption{Published PTB-XL superdiagnostic results for context. ``Operating
points'' is the number of embedding dimensionalities served by a single
trained artefact.}
\label{tab:sota}
\footnotesize
\setlength{\tabcolsep}{4pt}
\begin{tabular}{lccc}
\toprule
Method & Macro AUC & Params & Operating pts. \\
\midrule
XResNet1D-101~\cite{strodthoff2021} & 0.925 & 34\,M & 1 \\
Inception1D~\cite{strodthoff2021}   & 0.921 & 4.2\,M & 1 \\
Attention CNN~\cite{yao2020}        & 0.930 & $\sim$25\,M & 1 \\
Transformer~\cite{yan2019}          & 0.918 & $\sim$20\,M & 1 \\
Distillation~\cite{li2023kd}        & 0.910 & $\sim$5\,M & 1 \\
TinyML ECG~\cite{sakr2023}          & 0.870 & $<$1\,M & 1 \\
\midrule
Inception1D-MRL (ours) & \R{IncAUCsmall}--\R{IncAUClarge} & \R{IncParams} & 6 \\
XResNet101-MRL (ours)  & \R{XResAUCsmall}--\R{XResAUClarge} & \R{XResParams} & 6 \\
\bottomrule
\end{tabular}
\end{table}
